\documentclass[12pt]{article}
\usepackage[T1]{fontenc}
\usepackage[utf8]{inputenc}
\usepackage{lmodern}
\usepackage{textcomp}
\usepackage{enumitem}
\usepackage{geometry}
\geometry{margin=1in}
\begin{document}
\begin{center}
Answer Key - History - K-12 Level Assessment 15 \
\small (Answers are ordered exactly as in the question paper.)
\end{center}

\section*{Multiple Choice}
\begin{enumerate}[label=\arabic*.]
\item \textbf{Answer:} C
\\ \textit{Options:} A) They were specialized for an environment that suddenly changed.; B) Their tool technology was less advanced.; C) They were unsuited to the hotter climates of the sub-Saharan savanna.; D) The Neandertal body type required greater amounts of calories to sustain.
\\ \textit{Note:} The statement is correct because Neandertals primarily inhabited cold environments in Europe and Asia, and there is no evidence to suggest they ever faced the hotter climates of the sub-Saharan savanna, which would not have directly contributed to their extinction.
\item \textbf{Answer:} A
\\ \textit{Options:} A) tribute in the form of gold, jade, feathers, cloth, and jewels.; B) oppressive social and religious control based on military conquest.; C) ceramic production, a distinctive architectural style, and artistic expression.; D) religious beliefs that required extensive and escalating human sacrifice.
\\ \textit{Note:} The Aztec Empire relied on a tribute system where conquered regions were required to provide valuable resources such as gold, jade, feathers, cloth, and jewels, which fueled the empire's wealth and power.
\item \textbf{Answer:} D
\\ \textit{Options:} A) careful planning and scheduling.; B) reliance on steady sources of food such as shellfish.; C) increased use of imported or "exotic" materials; D) little planning
\\ \textit{Note:} Opportunistic foraging is characterized by individuals taking advantage of readily available food sources without extensive planning or strategy, reflecting a more spontaneous approach to obtaining sustenance.
\item \textbf{Answer:} B
\\ \textit{Options:} A) happened everywhere except in South America.; B) took thousands of years.; C) took hundreds of thousands of years.; D) always required careful planning.
\\ \textit{Note:} The transition from food gathering to a total reliance on agriculture took thousands of years because it involved gradual changes in human society, including the domestication of plants and animals, shifts in settlement patterns, and cultural adaptations that required time to develop and stabilize.
\item \textbf{Answer:} B
\\ \textit{Options:} A) biological anthropology.; B) anthropology.; C) cultural anthropology.; D) paleoanthropology.
\\ \textit{Note:} Archaeology is considered a subfield of anthropology because it focuses on the study of human cultures and societies through their material remains, which is a central aspect of anthropological research.
\end{enumerate}

\section*{True / False}
\begin{enumerate}[label=\arabic*.]
\setcounter{enumi}{5}
\item \textbf{Answer:} True
\\ \textit{Options:} A) True; B) False
\\ \textit{Note:} The average cranial capacity of Homo erectus is estimated to be around 900 to 1,100 cubic centimeters, which supports the statement that it is just under 1000 cc.
\item \textbf{Answer:} False
\\ \textit{Options:} A) True; B) False
\\ \textit{Note:} The statement is false because while Homo erectus had a larger brain size compared to earlier hominins, the development of art and aesthetic sense is more closely associated with later species, such as Homo sapiens, who demonstrated clear evidence of artistic expression and symbolic thought.
\item \textbf{Answer:} True
\\ \textit{Options:} A) True; B) False
\\ \textit{Note:} Teosinte is considered the wild ancestor of maize because genetic studies have shown that maize was domesticated from teosinte's species, which provided the necessary traits for the development of modern corn.
\item \textbf{Answer:} False
\\ \textit{Options:} A) True; B) False
\\ \textit{Note:} The Gault site is located inland in Central Texas, far from navigable rivers and the Gulf coast, which contradicts the claim that Paleoindians selected it for proximity to such water sources.
\item \textbf{Answer:} True
\\ \textit{Options:} A) True; B) False
\\ \textit{Note:} Stonehenge is indeed a megalithic monument in England characterized by its iconic arrangement of large stones in a circular formation, which is a defining feature of its historical significance.
\end{enumerate}

\section*{Long Form Response}
\begin{enumerate}[label=\arabic*.]
\setcounter{enumi}{10}
\item \textbf{Answer:} Catalhoyuk is classified as a large village rather than a city due to several key features, particularly in its architecture and social structure. Unlike cities, which typically have palaces and public buildings, Catalhoyuk lacked such monumental architecture, indicating a more egalitarian society. Additionally, there is little evidence of status differentiation among its inhabitants, suggesting that social hierarchies were not pronounced. The compact, interconnected houses of the village further reflect a communal way of life rather than the individualism often found in urban centers. These characteristics highlight Catalhoyuk's identity as a large village where community and cooperation were emphasized over the complexities of city life.
\\ \textit{Note:} The answer correctly identifies that Catalhoyuk's lack of monumental architecture and pronounced social hierarchies, along with its compact and interconnected housing structure, collectively indicate its classification as a large village rather than a city.
\item \textbf{Answer:} Humans share a common ancestry with primates, which means that we both evolved from a shared ancestor millions of years ago. This evolutionary relationship is significant because it helps us understand the traits we have in common with other primates, such as social behaviors and cognitive abilities. However, there are key differences that define the human species, including our ability to use complex language, create intricate tools, and develop advanced cultures. These differences have allowed humans to adapt and thrive in diverse environments, setting us apart from our primate relatives. Understanding both our similarities and differences with primates helps us appreciate the unique traits that define humanity.
\\ \textit{Note:} The answer correctly emphasizes the shared ancestry between humans and primates, illustrating how this common evolutionary background informs our understanding of both the similarities and the unique traits that distinguish humans, such as advanced language, tool use, and cultural development.
\end{enumerate}

\vfill
\noindent\textit{End of Answer Key}
\end{document}